\documentclass[11pt,a4paper]{article}

\usepackage[utf8]{inputenc}
\usepackage[T1]{fontenc}
\usepackage{lmodern}
\usepackage{microtype}
\usepackage{amsmath,amssymb,bm}
\usepackage{graphicx}
\usepackage{booktabs}
\usepackage{hyperref}
\usepackage{geometry}
\usepackage{xcolor}
\usepackage{listings}
\usepackage{caption}
\usepackage{subcaption}
\usepackage{siunitx}
\geometry{margin=2.7cm}

\hypersetup{
    colorlinks=true,
    linkcolor=blue!45!black,
    citecolor=blue!45!black,
    urlcolor=blue!45!black
}

\lstset{
    basicstyle=\ttfamily\small,
    breaklines=true,
    frame=single,
    columns=fullflexible
}

\title{FYS5419 Project 2\\\large Report Template}
\author{Your Name\\University of Oslo}
\date{Spring 2026}

\begin{document}
\maketitle

\begin{abstract}
Write a concise summary of the problem, method, main numerical results, and conclusion. Replace this template text after you have chosen one of the Project 2 alternatives.
\end{abstract}

\tableofcontents
\newpage

\section{Introduction}
Explain the scientific problem and why the selected quantum algorithm or quantum machine learning model is relevant. State clearly which Project 2 alternative you chose.

\paragraph{Chosen path.} TODO: Select exactly one of the following paths and remove the unused material from the final report:
\begin{enumerate}
    \item Quantum Fourier Transform and Quantum Phase Estimation.
    \item Quantum Machine Learning with the Iris data set.
    \item Variational Quantum Boltzmann Machines.
    \item Adaptive VQE / eigenvalue problems.
\end{enumerate}

\section{Theory}

\subsection{Quantum Fourier Transform path}
For the QFT path, start from the definition
\begin{equation}
    |j\rangle \mapsto \frac{1}{\sqrt{N}}\sum_{k=0}^{N-1} e^{2\pi i j k/N}|k\rangle,
\end{equation}
and show that the transformation is unitary. Include the one-, two-, and three-qubit matrices, then derive the gate decomposition and inverse QFT.

\subsection{Quantum Phase Estimation path}
Describe the eigenvalue equation
\begin{equation}
    U|u\rangle = e^{2\pi i \phi}|u\rangle
\end{equation}
and how the inverse QFT extracts the phase estimate \(\tilde{\phi}\). Discuss the relation between phase and energy for a time-evolution operator \(U=\exp(-iHt)\).

\subsection{Quantum Machine Learning path}
For the QML path, describe the feature map, ansatz, measurement rule, loss function, and parameter-shift gradient. For binary classification, define the prediction \(f(\bm{x};\bm{\theta})\) and cross-entropy loss.

\section{Numerical methods and implementation}
Describe the code structure, algorithms, and validation tests. Put reusable code in \texttt{src/fys5419\_project\_2/}; scripts used to generate results can go in \texttt{scripts/}. Include pseudocode where useful.

\subsection{Reproducibility}
Record the Python version, package versions, random seeds, and all parameters used to generate the final figures and tables.

\section{Results}
Place final figures in \texttt{results/figures/} and final tables in \texttt{results/tables/}. Every table and figure should have a label, caption, and clear axes/units where applicable.

\begin{table}[h]
    \centering
    \caption{Example results table. Replace with your final numerical results.}
    \label{tab:example}
    \begin{tabular}{lcc}
        \toprule
        Method & Error & Runtime \\
        \midrule
        Reference & TODO & TODO \\
        Project implementation & TODO & TODO \\
        \bottomrule
    \end{tabular}
\end{table}

\section{Discussion}
Interpret the results. Discuss accuracy, numerical stability, scaling, and limitations. For comparisons, explain what was held fixed and what was varied.

\section{Conclusion}
Summarize what was implemented, what was validated, and what you learned. Mention possible improvements or extensions.

\section*{Acknowledgements and collaboration}
State collaborators, division of work, and any AI/tool assistance according to the course requirements.

\appendix
\section{Code usage}
Give the exact commands used to reproduce the final results. Example:

\begin{lstlisting}[language=bash]
python scripts/run_qft_reference.py
make pdf
\end{lstlisting}

\section{Additional derivations}
Use the appendix for derivations or extra results that would interrupt the main text.

\begin{thebibliography}{9}
\bibitem{hundt2022}
Robert Hundt, \emph{Quantum Computing for Programmers}, Cambridge University Press, 2022.

\bibitem{nielsenchuang2000}
Michael A. Nielsen and Isaac L. Chuang, \emph{Quantum Computation and Quantum Information}, Cambridge University Press, 2000.

\bibitem{coursegithub}
CompPhysics, \emph{QuantumComputingMachineLearning course repository}, \url{https://github.com/CompPhysics/QuantumComputingMachineLearning}.
\end{thebibliography}

\end{document}
